% =====================================================================
%  preamble.tex —— Kollár・森《双有理几何》中文译注本 · 导言区
%  Lead 独占；译者只读，不得修改
% =====================================================================
\usepackage[b5paper,top=21mm,bottom=21mm,left=19mm,right=19mm,headsep=5mm,footskip=9mm]{geometry}
\usepackage{amsmath,amssymb,mathrsfs,bm}
\usepackage{graphicx}
\usepackage{xcolor}
\usepackage{booktabs,array,multirow,longtable,multicol}
\usepackage{newunicodechar}% 使正文偶发的字面 ■/□ 也能正确排版
\newunicodechar{■}{\ensuremath{\blacksquare}}
\newunicodechar{□}{\ensuremath{\square}}
% 译者偶发的 Unicode 数学符号（应写在 $...$ 内，这里作兜底，避免编译缺字）
\newunicodechar{≥}{\ensuremath{\geq}}
\newunicodechar{≤}{\ensuremath{\leq}}
\newunicodechar{≠}{\ensuremath{\neq}}
\newunicodechar{≡}{\ensuremath{\equiv}}
\newunicodechar{≅}{\ensuremath{\cong}}
\newunicodechar{∼}{\ensuremath{\sim}}
\newunicodechar{≈}{\ensuremath{\approx}}
\newunicodechar{±}{\ensuremath{\pm}}
\newunicodechar{×}{\ensuremath{\times}}
\newunicodechar{⋅}{\ensuremath{\cdot}}
\newunicodechar{⊗}{\ensuremath{\otimes}}
\newunicodechar{⊕}{\ensuremath{\oplus}}
\newunicodechar{⌈}{\ensuremath{\lceil}}
\newunicodechar{⌉}{\ensuremath{\rceil}}
\newunicodechar{⌊}{\ensuremath{\lfloor}}
\newunicodechar{⌋}{\ensuremath{\rfloor}}
\newunicodechar{∞}{\ensuremath{\infty}}
\newunicodechar{→}{\ensuremath{\to}}
\newunicodechar{↦}{\ensuremath{\mapsto}}
\newunicodechar{⊂}{\ensuremath{\subset}}
\newunicodechar{⊆}{\ensuremath{\subseteq}}
\newunicodechar{∈}{\ensuremath{\in}}
\newunicodechar{Δ}{\ensuremath{\Delta}}
\newunicodechar{Σ}{\ensuremath{\Sigma}}
\newunicodechar{Γ}{\ensuremath{\Gamma}}
\newunicodechar{Λ}{\ensuremath{\Lambda}}
\newunicodechar{ξ}{\ensuremath{\xi}}
\newunicodechar{η}{\ensuremath{\eta}}
\newunicodechar{λ}{\ensuremath{\lambda}}
\newunicodechar{μ}{\ensuremath{\mu}}
\newunicodechar{ν}{\ensuremath{\nu}}
\newunicodechar{ω}{\ensuremath{\omega}}
\newunicodechar{φ}{\ensuremath{\varphi}}
\newunicodechar{ψ}{\ensuremath{\psi}}
\newunicodechar{・}{\textperiodcentered}
\usepackage{enumitem}
\usepackage{fancyhdr}
\usepackage{caption}
\usepackage{tcolorbox}
\tcbuselibrary{most}
\usepackage{tikz}
\usepackage{tikz-cd}
\usepackage{hyperref}
\hypersetup{colorlinks=true,linkcolor=black,citecolor=black,urlcolor=blue!45!black,
            bookmarksnumbered=true,pdftitle={双有理几何（Kollár・森重文 著，中文译注本）},
            pdfauthor={János Kollár, 森重文}}

\setlength{\parindent}{2em}
\linespread{1.12}
\allowdisplaybreaks[2]
\setlength{\emergencystretch}{6em}
\tolerance=3000
\hbadness=10000

% ---------- 页眉页脚 ----------
\pagestyle{fancy}
\fancyhf{}
\renewcommand{\headrulewidth}{0.3pt}
\fancyhead[LE]{\small\thepage}
\fancyhead[RE]{\small\leftmark}
\fancyhead[RO]{\small\thepage}
\fancyhead[LO]{\small\rightmark}
\fancyfoot{}
\renewcommand{\chaptermark}[1]{\markboth{\CTEXthechapter\quad #1}{}}
\renewcommand{\sectionmark}[1]{\markright{\thesection\quad #1}}

% ---------- 章节标题格式 ----------
\ctexset{
  chapter = {
    format = \centering\Large\bfseries,
    name = {第,章},
    number = \arabic{chapter},
    beforeskip = 18pt,
    afterskip = 22pt,
  },
  section = {
    format = \large\bfseries,
    beforeskip = 14pt,
    afterskip = 10pt,
  },
  subsection = {
    format = \normalsize\bfseries,
    beforeskip = 11pt,
    afterskip = 6pt,
  },
}

% ---------- 公式编号：章.序号（原书同此制；实际一律 \tag 手工给号）----------
\numberwithin{equation}{chapter}

% ---------- 陈述环境 ----------
% 本书记号规则（原书“绪论”第 (16) 条）：证明末尾 ■；定义、例等末尾 □；
% 定理、命题等：不给证明时 ■，证明在后面给出时 □。
% 故：定理类默认 ■，星号版 □；定义/例类默认 □，星号版 ■。
% 用法：\begin{theorem}{1.3} 正文…… \end{theorem}
% 可选参数 [附加名] 用于“定理 0.2（广中定理）”这种写法；编号仍为必选的第一参数。
\newcommand{\fkstmt}[3][]{\par\medskip\noindent\textbf{#2~#3}%
  \if\relax\detokenize{#1}\relax\else\textbf{（#1）}\fi\ \ignorespaces}
\newenvironment{definition}[2][]{\fkstmt[#1]{定义}{#2}}{\hfill$\square$\par\medskip}
\newenvironment{definition*}[2][]{\fkstmt[#1]{定义}{#2}}{\hfill$\blacksquare$\par\medskip}
\newenvironment{theorem}[2][]{\fkstmt[#1]{定理}{#2}}{\hfill$\blacksquare$\par\medskip}
\newenvironment{theorem*}[2][]{\fkstmt[#1]{定理}{#2}}{\hfill$\square$\par\medskip}
\newenvironment{proposition}[2][]{\fkstmt[#1]{命题}{#2}}{\hfill$\blacksquare$\par\medskip}
\newenvironment{proposition*}[2][]{\fkstmt[#1]{命题}{#2}}{\hfill$\square$\par\medskip}
\newenvironment{lemma}[2][]{\fkstmt[#1]{引理}{#2}}{\hfill$\blacksquare$\par\medskip}
\newenvironment{lemma*}[2][]{\fkstmt[#1]{引理}{#2}}{\hfill$\square$\par\medskip}
\newenvironment{corollary}[2][]{\fkstmt[#1]{推论}{#2}}{\hfill$\blacksquare$\par\medskip}
\newenvironment{corollary*}[2][]{\fkstmt[#1]{推论}{#2}}{\hfill$\square$\par\medskip}
\newenvironment{example}[2][]{\fkstmt[#1]{例}{#2}}{\hfill$\square$\par\medskip}
\newenvironment{example*}[2][]{\fkstmt[#1]{例}{#2}}{\hfill$\blacksquare$\par\medskip}
\newenvironment{remark}[2][]{\fkstmt[#1]{注}{#2}}{\par\medskip}
\newenvironment{principle}[2][]{\fkstmt[#1]{原理}{#2}}{\hfill$\square$\par\medskip}
\newenvironment{principle*}[2][]{\fkstmt[#1]{原理}{#2}}{\hfill$\blacksquare$\par\medskip}
\newenvironment{notation}[2][]{\fkstmt[#1]{记号}{#2}}{\par\medskip}
\newenvironment{question}[2][]{\fkstmt[#1]{问}{#2}}{\hfill$\blacksquare$\par\medskip}
\newenvironment{problem}[2][]{\fkstmt[#1]{问题}{#2}}{\hfill$\blacksquare$\par\medskip}
\newenvironment{conjecture}[2][]{\fkstmt[#1]{猜想}{#2}}{\hfill$\blacksquare$\par\medskip}
\newenvironment{conjecture*}[2][]{\fkstmt[#1]{猜想}{#2}}{\hfill$\square$\par\medskip}
\newenvironment{problem*}[2][]{\fkstmt[#1]{问题}{#2}}{\hfill$\square$\par\medskip}
\newenvironment{question*}[2][]{\fkstmt[#1]{问}{#2}}{\hfill$\square$\par\medskip}
\newenvironment{counterexample}[2][]{\fkstmt[#1]{反例}{#2}}{\hfill$\blacksquare$\par\medskip}
\newenvironment{assumption}[2][]{\fkstmt[#1]{假设}{#2}}{\hfill$\square$\par\medskip}
\newenvironment{formulaenv}[2][]{\fkstmt[#1]{公式}{#2}}{\par\medskip}

% ---------- 第三种尾符情形：证明紧跟其后（原书声明末尾**无符号**，■ 在证明末）----------
% 约定：无后缀=书中不证（■／定义类 □）；* 号=证明在后文（□／定义类 ■）；后缀 p=证明紧跟（无符号）
\newenvironment{theoremp}[2][]{\fkstmt[#1]{定理}{#2}}{\par\medskip}
\newenvironment{propositionp}[2][]{\fkstmt[#1]{命题}{#2}}{\par\medskip}
\newenvironment{lemmap}[2][]{\fkstmt[#1]{引理}{#2}}{\par\medskip}
\newenvironment{corollaryp}[2][]{\fkstmt[#1]{推论}{#2}}{\par\medskip}
\newenvironment{definitionp}[2][]{\fkstmt[#1]{定义}{#2}}{\par\medskip}
\newenvironment{examplep}[2][]{\fkstmt[#1]{例}{#2}}{\par\medskip}
\newenvironment{principlep}[2][]{\fkstmt[#1]{原理}{#2}}{\par\medskip}
\newenvironment{questionp}[2][]{\fkstmt[#1]{问}{#2}}{\par\medskip}
\newenvironment{proof}[1][证明]{\par\medskip\noindent\textbf{#1}\ \ignorespaces}{\hfill$\blacksquare$\par\medskip}
\newcommand{\proofless}{}

% ---------- 图题（手工编号，与原书一致）----------
\newcommand{\figtitle}[2]{\par\smallskip\noindent{\centering\small 图~#1\quad #2\par}\smallskip}

% ---------- 表题 ----------
\newcommand{\tabtitle}[2]{\par\smallskip\noindent{\centering\small 表~#1\quad #2\par}\smallskip}

% ---------- “本节重点”与“章末评注”（全局开关）----------
\newif\ifshowtakeaway     \showtakeawaytrue
\newif\ifshowchapterreview \showchapterreviewtrue

\newtcolorbox{takeawaybox}{
  enhanced, breakable, colback=blue!4!white, colframe=blue!55!black,
  boxrule=0.4pt, arc=1.5pt, left=5pt, right=5pt, top=3pt, bottom=3pt,
  title={\small\bfseries 本节重点}, fonttitle=\small\bfseries,
  coltitle=white, colbacktitle=blue!55!black, attach boxed title to top left={xshift=4pt,yshift=-2pt},
}
\newtcolorbox{reviewbox}{
  enhanced, breakable, colback=orange!3!white, colframe=orange!65!black,
  boxrule=0.4pt, arc=1.5pt, left=6pt, right=6pt, top=4pt, bottom=4pt,
  title={\small\bfseries 本章回顾与评注}, fonttitle=\small\bfseries,
  coltitle=white, colbacktitle=orange!65!black, attach boxed title to top left={xshift=4pt,yshift=-2pt},
}
\newcommand{\takeaway}[1]{\ifshowtakeaway\begin{takeawaybox}\small #1\end{takeawaybox}\fi}
\newcommand{\chapterreview}[1]{\ifshowchapterreview\begin{reviewbox}\small #1\end{reviewbox}\fi}

% ---------- 常用记号 ----------
\newcommand{\R}{\mathbb{R}}
\newcommand{\C}{\mathbb{C}}
\newcommand{\Z}{\mathbb{Z}}
\newcommand{\Q}{\mathbb{Q}}
\newcommand{\N}{\mathbb{N}}
\newcommand{\PP}{\mathbb{P}}
\newcommand{\PPn}{\mathbb{P}^{n}}
\newcommand{\rd}{\mathrm{d}}
\newcommand{\id}{\mathrm{id}}
\newcommand{\cO}{\mathcal{O}}
\newcommand{\cF}{\mathcal{F}}
\newcommand{\cG}{\mathcal{G}}
\newcommand{\cH}{\mathcal{H}}
\newcommand{\cI}{\mathcal{I}}
\newcommand{\cJ}{\mathcal{J}}
\newcommand{\cK}{\mathcal{K}}
\newcommand{\cL}{\mathcal{L}}
\newcommand{\cM}{\mathcal{M}}
\newcommand{\cN}{\mathcal{N}}
\newcommand{\cP}{\mathcal{P}}
\newcommand{\cQ}{\mathcal{Q}}
\newcommand{\cR}{\mathcal{R}}
\newcommand{\cS}{\mathcal{S}}
\newcommand{\cT}{\mathcal{T}}
\newcommand{\cC}{\mathcal{C}}
\newcommand{\cD}{\mathcal{D}}
\newcommand{\cA}{\mathcal{A}}
\newcommand{\cB}{\mathcal{B}}
\newcommand{\cX}{\mathcal{X}}
\newcommand{\cZ}{\mathcal{Z}}
\newcommand{\cE}{\mathcal{E}}
\DeclareMathOperator{\Div}{Div}
\DeclareMathOperator{\Pic}{Pic}
\DeclareMathOperator{\Cl}{Cl}
\DeclareMathOperator{\Supp}{Supp}
\DeclareMathOperator{\Exc}{Exc}
\DeclareMathOperator{\Sing}{Sing}
\DeclareMathOperator{\Spec}{Spec}
\DeclareMathOperator{\Proj}{Proj}
\DeclareMathOperator{\Nef}{Nef}
\DeclareMathOperator{\Mov}{Mov}
\DeclareMathOperator{\Eff}{Eff}
\DeclareMathOperator{\mult}{mult}
\DeclareMathOperator{\ord}{ord}
\DeclareMathOperator{\rank}{rank}
\DeclareMathOperator{\im}{Im}
\DeclareMathOperator{\Ker}{Ker}
\DeclareMathOperator{\Coker}{Coker}
\newcommand{\KK}{K}
\newcommand{\klt}{\mathrm{klt}}
\newcommand{\lc}{\mathrm{lc}}
\newcommand{\NE}{\overline{NE}}
\newcommand{\vp}{\varphi}
\newcommand{\eps}{\varepsilon}
% 图内文本用：无衬线小号
\newcommand{\figfont}{\small}

% ---------- 参考文献（原书式：条目标号 [Abh56]）----------
\newenvironment{biblist}{\begin{list}{}{\leftmargin=3.4em\itemindent=-3.4em
  \itemsep=0.6pt\parsep=0pt\topsep=2pt}}{\end{list}}
\newcommand{\bibentry}[2]{\item[\textnormal{[#1]}] #2\par}

% ---------- 索引 / 术语表 ----------
\newenvironment{termlist}{\begin{list}{}{\leftmargin=2.2em\itemindent=0pt
  \itemsep=0.4pt\parsep=0pt\topsep=2pt}}{\end{list}}
\newcommand{\termitem}[2]{\item #1\hfill{\small #2}}
